\section*{Chapter 4: Curving Spacetime}
\begin{description}[style=multiline,leftmargin=2em,labelwidth=1.5em]

% DOCX paragraph 2032
\item[1.]\hypertarget{note:04-01}{}As you will remember from Chapter 2, the gravitational constant \emph{G} is equal to
6.67408\(\times\)10\textsuperscript{--11} m\textsuperscript{3}/(kg
s\textsuperscript{2}), where m, kg, and s stand respectively for metres, kilograms and
seconds, which are units in the International System of Units (SI). It follows that the
quantity \emph{GM/(c}\textsuperscript{2} \emph{R} ) does not have units, as they all
cancel out. Hence, the quantity \emph{GM/(c}\textsuperscript{2} \emph{R} ) is a
\emph{pure number}.

% DOCX paragraph 2033
\item[2.]\hypertarget{note:04-02}{}The symbol \(\oplus\) is used in astronomy to indicate the planet Earth; therefore, for
example, \(M_{\oplus}\) indicates the mass of our planet. Similarly, the symbol 8, which
we will use shortly, indicates the Sun and any quantity relative to it.

% DOCX paragraph 2034
\item[3.]\hypertarget{note:04-03}{}The parsec is a unit of length commonly used in astronomy as it allows us to account for
the enormous distances encountered outside the solar system. For this reason kilometres
are never used in astronomy, but rather light years and their multiples. I recall that a
``light year'' does not measure an interval of time but of space: the space covered by a
beam of light when travelling for one year. When expressed in kilometres, this
corresponds to 10,000 billion (10\textsuperscript{13}) kilometres. On the other hand, a
parsec, corresponds to 3.26 light years, that is, about 30,000 billion kilometres.
Clearly, it is much easier to talk of one parsec than tens of thousands of billions of
kilometres~\ldots{}

\end{description}
